% Figure: symmetric double-wedge (diamond) airfoil at angle of attack in
% supersonic flow, with the four surface panels and their wave systems.
% Compression turns (into the flow) run through oblique shocks; expansion
% turns (away from the flow) run through Prandtl-Meyer fans.
\begin{figure}[htbp]
\centering
\begin{tikzpicture}[scale=1.0]
% diamond section: nose at left, tail at right, mid at top and bottom
\coordinate (N) at (0,0);      % nose (leading edge)
\coordinate (T) at (6,0);      % tail (trailing edge)
\coordinate (Up) at (3,0.9);   % upper mid-chord
\coordinate (Lo) at (3,-0.9);  % lower mid-chord
\draw[engblack, very thick] (N) -- (Up) -- (T) -- (Lo) -- cycle;
% panel labels
\node[font=\scriptsize] at (1.4,0.65) {panel 1 (upper fore)};
\node[font=\scriptsize] at (4.6,0.65) {panel 2 (upper aft)};
\node[font=\scriptsize] at (1.4,-0.7) {panel 3 (lower fore)};
\node[font=\scriptsize] at (4.6,-0.7) {panel 4 (lower aft)};
% freestream arrow (slightly below horizontal to indicate positive alpha)
\draw[enggray, -{Latex[length=1.8mm]}] (-2.4,-0.35) -- (-0.7,-0.15);
\node[enggray, font=\scriptsize, anchor=south] at (-1.7,-0.15) {$M_\infty$, $\alpha$};
% leading-edge oblique shocks (upper and lower), red
\draw[engred, thick] (N) -- (-1.1,1.4);
\draw[engred, thick] (N) -- (-1.1,-1.4);
\node[engred, font=\scriptsize, anchor=east] at (-1.1,1.2) {L.E. shocks};
% mid-chord expansion fans (upper and lower), blue fan lines
\foreach \a in {30,45,60} {
  \draw[engblue] (Up) -- ++(\a:1.3);
}
\foreach \a in {-30,-45,-60} {
  \draw[engblue] (Lo) -- ++(\a:1.3);
}
\node[engblue, font=\scriptsize, anchor=south west] at (3.3,1.6) {expansion fans};
% trailing-edge shocks/recompression
\draw[engred, thick] (T) -- (7.1,1.1);
\draw[engred, thick] (T) -- (7.1,-1.1);
\node[engred, font=\scriptsize, anchor=west] at (7.0,0.9) {T.E. shocks};
\end{tikzpicture}
\caption[Double-wedge diamond airfoil in supersonic flow]{A symmetric
double-wedge (diamond) section at angle of attack $\alpha$ in a
supersonic stream. The four straight panels each meet the flow at a
definite angle, and the wave system follows from walking the surface
panel by panel. Each leading-edge face that turns into the flow runs
through an attached oblique shock (red); each mid-chord shoulder that
turns the flow away from itself runs through a Prandtl-Meyer expansion
fan (blue); the flows off the upper and lower aft panels recompress
through trailing-edge shocks to match pressure and direction in the
wake. The surface pressure on each panel is whatever the running tally
of compressions and expansions delivers, and the resultant of the four
panel pressures is the exact inviscid lift and wave drag. Unlike the
flat plate, the diamond carries a thickness (volume) wave drag that
survives even at zero lift, because the fore panels compress and the
aft panels expand whatever the angle of attack.}
\label{fig:vol03ch13:diamond-airfoil}
\end{figure}
